% Ejercicio 4. Unir ambos sonidos en uno nuevo para escuchar el nombre y apellido correctamente.
Con la función \texttt{pastew()} se unen dos sonidos. Hay que prestar atención al orden, pues el primer argumento será el segundo sonido que se escuche y el segundo argumento será el primer sonido que se escuche.
\begin{lstlisting}[caption={Función para unir dos sonidos}, label={lst:union_sonidos}]
union <- pastew(apellidos, nombre, output="Wave")
\end{lstlisting}

En la Figura \ref{fig:img4_1} se puede observar la información del sonido.
\begin{figure}[H]
    \centering
    \includegraphics[width=0.95\linewidth]{images/4_1union.png}
    \caption{Unión de dos sonidos}
    \label{fig:img4_1}
\end{figure}

En el Cuadro \ref{tab:tabla2} se pueden observar los datos de los tres ficheros. La duración del fichero \textit{union} es la suma de las duraciones del fichero \textit{nombre} y el fichero \textit{apellidos}, aproximadamente. Sucede lo mismo con el número de muestras: $32,879 + 49,007 = 81,886$.
\begin{table}[ht]
    \centering
    \caption{Comparativa de los objetos Wave: \texttt{nombre}, \texttt{apellidos} y \texttt{union}}
    \label{tab:tabla2}
    \begin{tabular}{llll}
        \toprule
        \textbf{Parámetro} & \textbf{Objeto: \texttt{nombre}} & \textbf{Objeto: \texttt{apellidos}} & \textbf{Objeto: \texttt{union}}\\
        \midrule
        Número de muestras & 32879 & 49007 & 81886\\
        Duración (segundos) & 0.68 & 1.02 & 1.71 \\
        Tasa de muestreo (Hz) & 48000 & 48000 & 48000\\
        Canales & Stereo & Stereo & Mono\\
        Formato PCM & TRUE & TRUE & TRUE\\
        Bit & 16 & 16 & 16\\
        \bottomrule
    \end{tabular}
\end{table}
